\item Un hidrómetro es un instrumento utilizado para determinar la densidad de los líquidos. En la figura P14.16 se muestra uno simple. El bulbo de una jeringa se presiona y libera para dejar que la atmósfera eleve una muestra del líquido de interés en un tubo que contiene una barra calibrada de densidad conocida. La barra, de longitud $L$ y densidad promedio $\rho_0$, flota parcialmente sumergida en el líquido de densidad $\rho$. Una longitud $h$ de la barra sobresale de la superficie del líquido. Demuestre que la densidad del líquido está dada por
$$ \rho = \frac{L \rho_0}{L - h} $$
